\section{Latar Belakang}

Perkembangan \textit{Internet of Things} (IoT)\label{first:iot} mendorong konektivitas perangkat, sensor, dan sistem dalam skala yang sangat besar.
Jumlah perangkat IoT diproyeksikan melampaui 60 miliar unit pada 2030 dengan karakteristik yang heterogen dari sisi kemampuan komputasi, energi, dan konteks penggunaan \parencite{ahsan_comprehensive_2025}.
Pertumbuhan ini meningkatkan tantangan keamanan, terutama pada autentikasi, integritas data, dan kontrol akses.
Keterbatasan sumber daya perangkat juga menyulitkan penerapan mekanisme keamanan konvensional secara langsung \parencite{dritsas_survey_2025}.

Kontrol akses menjadi salah satu komponen penting dalam keamanan IoT karena menentukan entitas yang dapat mengakses sumber daya pada kondisi tertentu.
Model konvensional seperti \textit{Role-Based Access Control} (RBAC)\label{first:rbac} dan \textit{Discretionary Access Control} (DAC)\label{first:dac} kurang sesuai untuk lingkungan IoT yang dinamis, berskala besar, dan kaya konteks \parencite{ragothaman_access_2023}.
Sebaliknya, \textit{Attribute-Based Access Control} (ABAC)\label{first:abac} memungkinkan keputusan akses berdasarkan kombinasi atribut subjek, objek, dan lingkungan sehingga lebih sesuai dengan heterogenitas dan dinamika IoT \parencite{ahsan_comprehensive_2025}.

Namun, perancangan kebijakan ABAC secara manual kompleks dan mahal sehingga organisasi sering mempertahankan mekanisme kontrol akses lama daripada menanggung biaya migrasi \parencite{diaz-rodriguez_abac_2025}.
\textit{Policy mining} kemudian dikembangkan untuk mengekstraksi kebijakan ABAC secara otomatis dari log akses dan data atribut.
Pendekatan metaheuristik telah digunakan untuk tujuan ini, antara lain \textit{Ant Colony Optimization} (ACO)\label{first:aco} untuk migrasi kebijakan lintas sistem heterogen \parencite{siyuan_abac_2025} dan \textit{Upgrade Black-Winged Kite} (UBK)\label{first:ubk} untuk optimasi kebijakan ABAC pada berbagai data \parencite{li_attribute_2026}.
Namun, keduanya ditujukan untuk pemrosesan \textit{batch/offline} pada \textit{server} dengan sumber daya memadai, bukan perangkat \textit{edge-IoT} yang terbatas.

Keterbatasan memori, komputasi, dan energi pada perangkat IoT menuntut mekanisme keamanan yang ringan agar dapat dijalankan langsung pada \textit{edge} tanpa bergantung sepenuhnya pada komputasi awan \parencite{canavese_security_2024}.
Pada optimasi metaheuristik untuk \textit{edge computing}, kebutuhan ini mulai dijawab melalui algoritma ringan dan adaptif, seperti \textit{Modified Greylag Goose Optimization} (MGGO)\label{first:mggo} untuk alokasi layanan IoT yang tidak memerlukan pelatihan atau pemodelan eksplisit yang berat \parencite{khosrowshahi_refined_2025}.

Prinsip tersebut belum banyak diterapkan pada \textit{policy mining} ABAC.
Metaheuristik seperti ACO dan UBK menggunakan pendekatan berbasis populasi sehingga tidak dirancang untuk eksekusi ringan dan berulang pada perangkat berdaya rendah.
Di sisi lain, pendekatan \textit{policy mining} inkremental yang lebih efisien masih didominasi metode berbasis aturan \parencite{batra_incremental_2021,bamberger_automated_2024}, sehingga tidak memanfaatkan kemampuan eksplorasi kombinatorial yang menjadi keunggulan metaheuristik.

Kesenjangan ini semakin terlihat pada \textit{edge gateway} di lingkungan seperti \textit{smart building} yang harus mengelola akses terhadap ratusan sensor, aktuator, dan data untuk berbagai pengguna dan perangkat.
Perancangan kebijakan secara manual bersifat mahal, lambat, subjektif, dan rentan kesalahan sehingga organisasi cenderung mempertahankan RBAC/DAC daripada melakukan migrasi \parencite{diaz-rodriguez_abac_2025}.
Padahal, kelemahan kontrol akses ini dapat menimbulkan dampak yang signifikan.
Pada studi IBM-Ponemon tahun 2024 yang dikutip oleh \textcite{diaz-rodriguez_abac_2025}, memperkirakan rata-rata biaya satu pelanggaran data dapat mencapai USD 4,88 juta.

\textit{Policy mining} dapat mengotomatisasi pembentukan kebijakan, tetapi masih menghadapi dua keterbatasan utama pada \textit{edge-IoT}.
Pertama, metode berbasis \textit{association rule} dan sejenisnya mengalami penurunan kinerja pada log berukuran besar serta peningkatan jumlah aturan ketika jumlah atribut bertambah \parencite{xu_mining_2015,inproceedings}.
Karena itu, metode mutakhir umumnya masih dirancang untuk \textit{server} dengan sumber daya yang memadai.
Kedua, \textit{policy mining} umumnya dilakukan secara \textit{batch}, sedangkan lingkungan IoT terus berubah akibat munculnya perangkat, peran, dan pola akses baru sehingga kebijakan perlu diperbarui secara berkala.
Pemrosesan di \textit{server} juga mengharuskan pengiriman log keluar perangkat, yang menambah beban jaringan, latensi, dan risiko privasi, sementara algoritma berbasis populasi dengan jejak memori besar sulit dijalankan pada \textit{edge gateway} terbatas.
Atas dasar tersebut, penelitian ini merancang algoritma \textit{policy mining} ABAC yang \textit{compact}, mempertimbangkan dependensi antar-atribut, dan mampu memperbarui kebijakan secara inkremental sehingga dapat dijalankan langsung pada perangkat \textit{edge-IoT} tanpa mengorbankan kualitas kebijakan.